\documentclass[11pt,letterpaper]{article}

% --- PAQUETES DE DISEÑO ---
\usepackage[utf8]{inputenc}
% es-nodecimaldot evita el conflicto de babel-español con siunitx
\usepackage[spanish, es-tabla, es-nodecimaldot]{babel}
\usepackage[T1]{fontenc}
\usepackage{lmodern}
\usepackage{microtype}
\usepackage{geometry}
\usepackage{graphicx}
\usepackage{float}
\usepackage{url}
\usepackage{hyperref}
\usepackage{xurl}          % corta las URLs largas para que no se salgan del margen
\usepackage{listings}
\usepackage{amsmath}
\usepackage{amssymb}
\usepackage{booktabs}
\usepackage{siunitx}
\usepackage{subcaption}
\usepackage[authoryear,round]{natbib}
\usepackage{tikz}
\usepackage{pgfplots}
\pgfplotsset{compat=1.18}
\usetikzlibrary{shapes.geometric, arrows.meta, positioning, calc}
\usepackage{adjustbox}
\geometry{top=2cm, bottom=2.5cm, left=2.2cm, right=2.2cm}
% Evita que identificadores largos en \texttt se salgan del margen.
\emergencystretch=2.5em

% --- COLORES ---
\usepackage{xcolor}
\definecolor{mainBlue}{RGB}{0, 32, 96}
\definecolor{lightGray}{RGB}{240, 240, 240}

% --- TÍTULOS Y SECCIONES ---
\usepackage{titlesec}
\titleformat{\section}{\large\bfseries\color{mainBlue}}{\thesection.}{0.5em}{}[\titlerule]

% --- CAJAS Y CÓDIGO ---
\usepackage{tcolorbox}

\lstset{
    basicstyle=\ttfamily\footnotesize,
    backgroundcolor=\color{lightGray!50},
    frame=single,
    rulecolor=\color{mainBlue!20},
    columns=fullflexible,
    breaklines=true,
    breakatwhitespace=true,
    showstringspaces=false,
    xleftmargin=10pt,
    xrightmargin=6pt,
    inputencoding=utf8,
    extendedchars=true,
    literate={á}{{\'a}}1 {é}{{\'e}}1 {í}{{\'i}}1 {ó}{{\'o}}1 {ú}{{\'u}}1 {ñ}{{\~n}}1,
}

\newtcolorbox{analisisBox}[1]{
    colback=lightGray!30, colframe=mainBlue, fonttitle=\bfseries,
    title=#1, arc=0pt, outer arc=0pt, width=\linewidth,
    left=4pt, right=4pt
}

% siunitx: separador de miles con espacio fino, coma decimal
\sisetup{group-separator={\,}, output-decimal-marker={,}, group-minimum-digits=4, group-digits=integer}

\hypersetup{colorlinks=true, linkcolor=mainBlue, urlcolor=mainBlue, citecolor=mainBlue,
            breaklinks=true}

% --- INICIO DEL DOCUMENTO ---
\begin{document}
% babel-spanish activa "<" y ">" como atajos para comillas latinas;
% eso choca con listings y con la sintaxis >={...} de TikZ.
\shorthandoff{<>"}

% --- ENCABEZADO COMPACTO ---
\begin{center}
    {\color{mainBlue}\hrule height 1.5pt}
    \vspace{0.4cm}
    {\Large \textbf{\texttt{Procesamiento Digital de Imágenes Satelitales y Datos Geoespaciales con Python}}}\\
    \vspace{0.4cm}
    Entrega Clasificación Supervisada\\
    \vspace{0.4cm}
    {\large \textit{Campus Virtual del Instituto de Altos Estudios Espaciales Mario Gulich}}\\
    \vspace{0.4cm}
    {\color{mainBlue!30}\hrule height 0.8pt}
    \vspace{0.2cm}

    \begin{tabular}{cc}
        \textbf{Autor 1} \\[0.2cm]
        Argenis Alexandra Daza Roa \\[0.1cm]
        \footnotesize \texttt{aadazar@udistrital.edu.co}
    \end{tabular}

    \vspace{0.5cm}
    \small \textbf{Fecha:} \today
    \vspace{0.2cm}
    {\color{mainBlue}\hrule height 1.5pt}
\end{center}

\begin{abstract}
\noindent
En este trabajo se clasificó la cobertura de la tierra de una escena Sentinel-2 L2A del 9 de enero de 2025
(tesela T19NBH) en el departamento de Arauca, Colombia. Se usaron las bandas B02, B03, B04, B08, B11 y B12
a 10~m y seis clases del primer nivel del sistema CIAF: construcciones, cultivos y parcelas, vegetación herbácea
descubierta, bosques y montes, cuerpos de agua y tierras eriales. Las muestras se digitalizaron en QGIS sobre
dos recortes de $5\times5$~km. Con ellas se entrenó un Random Forest, se buscó la profundidad máxima por
validación cruzada y, a modo de comparación, se probó un KNN. Sobre el recorte de entrenamiento la exactitud
global fue de \num{0.979} ($\kappa=\num{0.972}$). Al validar con el segundo recorte, que tiene muestras
independientes, bajó a \num{0.668} ($\kappa=\num{0.531}$). La diferencia se explica sobre todo por tres cosas:
los suelos descubiertos y brillantes (playones, parcelas aradas) terminan clasificados como construcción, los
cultivos se confunden con la vegetación herbácea, y en el segundo recorte hay nubes y sombras que no hacían
parte del entrenamiento. Al final se plantean los ajustes para una segunda clasificación.
\end{abstract}

\tableofcontents
\vspace{0.6cm}

% =====================================================================
\section{Introducción}
% =====================================================================

Conocer cómo está cubierto el territorio (dónde hay bosque, dónde se cultiva, por dónde corre el agua, qué
tanto ha crecido lo urbano) es la base para casi cualquier decisión de ordenamiento o de gestión ambiental.
Hacer ese levantamiento en campo es costoso y lento, por eso desde hace décadas se recurre a las imágenes
satelitales. Cada tipo de superficie refleja la energía de una manera distinta a lo largo del espectro
electromagnético, y esa \emph{firma espectral} es la que permite separarlas de forma automática
\citep{chuvieco2010}.

La clasificación supervisada aprovecha justamente eso: a partir de un conjunto de píxeles de los que ya se
conoce la clase (las muestras de entrenamiento), un algoritmo aprende a asignar una clase al resto de la
imagen \citep{gulich}. Uno de los algoritmos más usados en teledetección es Random Forest, que combina
muchos árboles de decisión y suele dar buenos resultados sin necesidad de normalizar los datos
\citep{breiman2001}.

Para este trabajo se eligió una escena Sentinel-2 sobre el departamento de Arauca, en la Orinoquía colombiana.
Es una zona de sabanas inundables, bosques de galería, ríos con meandros, potreros y parcelas agrícolas, con el
casco urbano de Arauca sobre el río del mismo nombre. Esa mezcla la hace interesante para clasificar, pero
también difícil, porque varias coberturas tienen respuestas espectrales parecidas. El flujo de trabajo sigue
las Notebooks 02c, 03a y 03b del curso: lectura de las bandas, generación de muestras, entrenamiento,
evaluación y, como segunda instancia de validación, la aplicación del modelo sobre un recorte distinto con
muestras propias (apartado ``Yapa'' de la Notebook 03b).

% =====================================================================
\section{Objetivos}
% =====================================================================
\subsection{Objetivo general}
Clasificar las coberturas de la tierra de dos sectores del departamento de Arauca a partir de una imagen
Sentinel-2 L2A, usando aprendizaje supervisado en Python, y evaluar qué tan bien se comporta el modelo
cuando se aplica a una zona distinta de aquella en la que fue entrenado.

\subsection{Objetivos específicos}
\begin{itemize}
    \item Leer las bandas B02, B03, B04, B08, B11 y B12 directamente desde el archivo \texttt{.SAFE.zip},
    remuestrear las de 20~m a 10~m y generar dos recortes de $5\times5$~km.
    \item Definir las clases a partir del primer nivel del sistema CIAF y digitalizar en QGIS las muestras
    de entrenamiento y de validación.
    \item Analizar la separabilidad de las clases mediante sus firmas espectrales y gráficos de dispersión.
    \item Entrenar un clasificador Random Forest, ajustar su profundidad máxima por validación cruzada y
    compararlo con KNN.
    \item Evaluar los modelos con la matriz de confusión, la precisión, el \emph{recall}, el F1, la
    exactitud global y el índice kappa, tanto en el recorte de entrenamiento como en el de validación.
    \item Generar los mapas clasificados como GeoTIFF y GeoPackage, e identificar las principales fuentes de
    error para plantear una segunda clasificación.
\end{itemize}

% =====================================================================
\section{Marco conceptual}
\label{sec:marco}
% =====================================================================

\subsection{Imágenes satelitales y firmas espectrales}
Una imagen satelital es una medición por teledetección: un instrumento a bordo del satélite registra la
energía reflejada o emitida por la superficie en determinadas regiones del espectro electromagnético, y
cada píxel queda localizado geográficamente \citep{gulich}. Los sensores multiespectrales no miden el
espectro completo, sino ``recortes'' de él, que se conocen como bandas espectrales.

La \textbf{firma espectral} es la variación de la reflectancia de un material en función de la longitud de
onda \citep{chuvieco2010}. Las firmas de las coberturas que aparecen en este trabajo tienen rasgos
bastante conocidos \citep{gulich}:

\begin{itemize}
    \item \textbf{Vegetación sana:} refleja poco en el visible, salvo un pequeño pico en el verde, porque la
    clorofila absorbe el azul y el rojo. En el infrarrojo cercano (NIR) refleja mucho por la estructura
    interna de las hojas. En el infrarrojo de onda corta (SWIR) la reflectancia depende sobre todo del
    contenido de agua: la vegetación seca refleja más.
    \item \textbf{Agua:} el agua limpia refleja algo en el azul y absorbe con fuerza en el NIR y el SWIR,
    por eso se ve oscura en esas bandas. Cuando trae sedimentos o materia orgánica, como los ríos de la
    Orinoquía, su reflectancia en el visible sube y la diferencia con el suelo húmedo se hace menor.
    \item \textbf{Suelo:} en general su reflectancia aumenta con la longitud de onda, aunque depende de la
    humedad, la materia orgánica, la textura y los minerales. El suelo seco y arenoso refleja mucho en todo
    el espectro, mientras que el suelo húmedo absorbe más en el NIR y el SWIR.
\end{itemize}

Para la interpretación visual se usan composiciones RGB. Además del color natural (B04-B03-B02), en este
trabajo fue muy útil el \textbf{falso color de vegetación} (B08-B04-B03). Ahí la vegetación vigorosa
aparece en rojo, el agua en tonos oscuros o azules, y las superficies construidas y los suelos descubiertos
en tonos cian o grisáceos.

\subsection{La misión Sentinel-2}
Sentinel-2 es una misión del programa Copernicus de la Agencia Espacial Europea (ESA). Lleva el instrumento
multiespectral MSI, con 13 bandas entre el visible y el SWIR y resoluciones de 10, 20 y 60~m
\citep{drusch2012,esa2015}. El producto \textbf{L2A} viene corregido atmosféricamente y expresa los
valores en reflectancia de superficie (BOA), por lo que es el más adecuado para clasificar. Se distribuye
en teselas de $100\times100$~km y formato JPEG2000. Dentro de la carpeta \texttt{GRANULE/.../IMG\_DATA}
las bandas se agrupan en \texttt{R10m}, \texttt{R20m} y \texttt{R60m} según su resolución.

\begin{table}[H]
\centering
\caption{Bandas Sentinel-2 usadas en la clasificación y su posición en el conjunto de atributos.}
\label{tab:bandas}
\small
\begin{tabular}{@{}llccc@{}}
\toprule
Banda & Región & Longitud de onda central (nm) & Resolución original (m) & Atributo \\
\midrule
B02 & Azul      & 490  & 10 & \texttt{Band\_0} \\
B03 & Verde     & 560  & 10 & \texttt{Band\_1} \\
B04 & Rojo      & 665  & 10 & \texttt{Band\_2} \\
B08 & NIR       & 842  & 10 & \texttt{Band\_3} \\
B11 & SWIR 1    & 1610 & 20 & \texttt{Band\_4} \\
B12 & SWIR 2    & 2190 & 20 & \texttt{Band\_5} \\
\bottomrule
\end{tabular}
\end{table}

\subsection{Clasificación supervisada}
En el aprendizaje supervisado se cuenta con un conjunto de valores $X$ y una variable de respuesta $Y$, y
se busca aprender una función $f(X)$ que permita predecir $Y$. Cuando $Y$ es una variable categórica se habla
de un problema de \textbf{clasificación} \citep{gulich}. En imágenes, $X$ son las reflectancias de cada
píxel en las distintas bandas, que definen el \emph{espacio de atributos}, y $Y$ es la clase de cobertura.
Cuanto más separadas estén las nubes de puntos de cada clase en ese espacio, más fácil será clasificarlas.

\paragraph{Árboles de decisión.} Son métodos no paramétricos que dividen los datos con preguntas del tipo
SI/NO sobre las variables, por ejemplo ``¿la reflectancia en el NIR es mayor que cierto umbral?''. En cada
nodo se elige la división que más reduce la incertidumbre, medida con el índice de Gini o con la entropía.
Un árbol sin límites termina ajustándose de más a los datos de entrenamiento (\emph{overfitting}), por eso
se controlan parámetros como la profundidad máxima (\texttt{max\_depth}) \citep{gulich}.

\paragraph{Random Forest.} Es un método de \emph{ensemble} basado en \emph{bagging}. Se entrenan muchos
árboles, cada uno con un subconjunto aleatorio de los datos y de los atributos, y la clase final es la más
votada \citep{breiman2001}. Es más robusto que un árbol solo y no requiere normalizar los datos. Como
contrapartida, cuesta más interpretarlo y no funciona bien con conjuntos de datos muy pequeños
\citep{gulich}. Sus parámetros principales son la cantidad de árboles (\texttt{n\_estimators}) y la
profundidad máxima.

\paragraph{KNN.} Asigna a cada píxel la clase mayoritaria entre sus $k$ vecinos más cercanos en el espacio de
atributos. Como trabaja con distancias, es sensible a la escala de las variables, aunque en bandas de un
mismo sensor las reflectancias están en rangos parecidos. En el curso se recomienda usar
\texttt{RobustScaler} o \texttt{MinMaxScaler} cuando se quiere escalar.

\subsection{Evaluación de la exactitud}
La \textbf{matriz de confusión} cruza la clase real con la predicha para cada muestra de prueba
\citep{congalton1991}. A partir de ella se calculan, para cada clase $i$:
\begin{equation}
\text{Precisión}(i)=\frac{TP_i}{TP_i+FP_i},\qquad
\text{Recall}(i)=\frac{TP_i}{TP_i+FN_i},\qquad
F1(i)=2\,\frac{\text{Precisión}(i)\cdot\text{Recall}(i)}{\text{Precisión}(i)+\text{Recall}(i)}
\end{equation}
En teledetección, $1-\text{Precisión}$ corresponde al \textbf{error de comisión} y $1-\text{Recall}$ al
\textbf{error de omisión}. La \textbf{exactitud global} es la proporción de aciertos sobre el total, y el
\textbf{índice kappa} descuenta el acuerdo que se obtendría por azar \citep{cohen1960}:
\begin{equation}
\kappa=\frac{p_o-p_e}{1-p_e}
\end{equation}
donde $p_o$ es la exactitud observada y $p_e$ la esperada por azar.

La validación cruzada repite varias veces la partición entre entrenamiento y prueba y promedia los
resultados, lo que da una estimación más estable para elegir hiperparámetros. Hay un detalle importante:
cuando los píxeles de prueba salen de los mismos polígonos que los de entrenamiento, son muy parecidos entre
sí (están autocorrelacionados), y las métricas tienden a ser optimistas. Por eso es tan útil una segunda
validación con un recorte distinto y muestras independientes.

\subsection{Sistema de clasificación CIAF}
Para definir las clases se tomó como referencia el sistema de categorías del Centro de Investigación en
Percepción Remota (CIAF), que organiza la cobertura y el uso de la tierra en cuatro niveles de detalle
\citep{ciaf}. Se trabajó con el primer nivel (exploratorio), que tiene seis categorías y es coherente con lo
que se puede distinguir en una imagen de 10~m sin trabajo de campo.

% =====================================================================
\section{Área de estudio y datos}
\label{sec:datos}
% =====================================================================

\subsection{Imagen Sentinel-2}
La imagen se descargó del Copernicus Browser \citep{copernicus}, filtrando por productos Sentinel-2 L2A:
\begin{center}
\small \url{https://browser.dataspace.copernicus.eu/?zoom=9&lat=7.2535&lng=-70.83984&themeId=DEFAULT-THEME&visualizationUrl=U2FsdGVkX19Bm\%2FbBYkLNpBX9pPZuAU1EZTmibP5ORrWV76CQCauUWJ5zDuKukkfUZO70\%2FuZejP4VpSQgBXFYHl18\%2BXOfRjUAsMpW8aOgYo8cukIgA7rZD6oIG6C\%2B6Y2w&datasetId=S2_L2A_CDAS}
\end{center}

\begin{table}[H]
\centering
\caption{Datos de la escena utilizada.}
\label{tab:escena}
\small
\begin{tabular}{@{}ll@{}}
\toprule
Producto & \texttt{S2A\_MSIL2A\_20250109T151721\_N0511\_R125\_T19NBH\_20250109T183248} \\
Satélite & Sentinel-2A \\
Nivel de procesamiento & L2A (reflectancia de superficie) \\
Fecha de adquisición & 9 de enero de 2025 \\
Órbita relativa / tesela & R125 / T19NBH \\
Sistema de referencia & WGS 84 / UTM zona 19N (EPSG:32619) \\
\bottomrule
\end{tabular}
\end{table}

La fecha cae en la temporada seca de la Orinoquía (diciembre a marzo). Esto ayuda a que haya pocas nubes,
pero también hace que haya más suelo descubierto: playones del río, parcelas preparadas o cosechadas y
pastos secos. Como se verá en los resultados, eso influyó bastante en la clasificación.

% COLOCAR FIGURA (opcional): pantallazo del Copernicus Browser con la escena y la zona seleccionada.

\subsection{Recortes}
Se trabajó con dos recortes cuadrados de $5\times5$~km ($500\times500$ píxeles), dibujados en QGIS como
polígonos (\texttt{recorte1.gpkg} y \texttt{recorte2.gpkg}). Están separados unos 80~km entre sí y tienen
paisajes distintos a propósito:

\begin{itemize}
    \item \textbf{Recorte 1 (entrenamiento):} centrado aproximadamente en $7{,}077^\circ$N,
    $70{,}764^\circ$W. Incluye el casco urbano de Arauca, el aeropuerto, el río con sus meandros, lagunas,
    bosques y algo de potreros. Hay pocas parcelas agrícolas.
    \item \textbf{Recorte 2 (validación):} centrado aproximadamente en $6{,}694^\circ$N,
    $71{,}381^\circ$W. Es un sector rural alrededor de un río, dominado por parcelas, sin centro poblado
    y con algunas nubes y sombras.
    % VERIFICAR: nombre del municipio y del río del recorte 2.
\end{itemize}

\begin{figure}[H]
\centering
\begin{subfigure}{0.48\linewidth}
    \centering
    \includegraphics[width=\linewidth]{vista_recorte1.png}
    \caption{Recorte 1 (entrenamiento).}
\end{subfigure}\hfill
\begin{subfigure}{0.48\linewidth}
    \centering
    \includegraphics[width=\linewidth]{vista_recorte2.png}
    \caption{Recorte 2 (validación).}
\end{subfigure}
% COLOCAR FIGURA: las dos vistas en color natural de los recortes, que están en la página 3 del PDF
% "entregaclasificacion - Colab" (también en salidas/vista_recorte1.png y salidas/vista_recorte2.png).
\caption{Recortes en color natural (B04-B03-B02).}
\label{fig:recortes}
\end{figure}

% =====================================================================
\section{Metodología}
\label{sec:metodologia}
% =====================================================================

Todo el procesamiento se hizo en Google Colab con un único script, \texttt{entregaclasificacion.py},
armado a partir de las funciones de las Notebooks 02c, 03a y 03b. La Figura~\ref{fig:flujo} resume el flujo.

\begin{figure}[H]
\centering
\begin{adjustbox}{max width=\linewidth}
\begin{tikzpicture}[
    node distance=0.55cm and 0.6cm,
    paso/.style={rectangle, rounded corners=2pt, draw=mainBlue, fill=lightGray!60, thick,
                 text width=3.1cm, align=center, minimum height=1.15cm, font=\footnotesize},
    flecha/.style={-{Stealth[length=2.2mm]}, thick, mainBlue}
]
\node[paso] (a) {Descarga escena S2 L2A (Copernicus Browser)};
\node[paso, right=of a] (b) {Lectura de bandas desde el \texttt{.zip} (02c)};
\node[paso, right=of b] (c) {Recortes y remuestreo B11, B12 a 10~m};
\node[paso, right=of c] (d) {Digitalización de muestras en QGIS};
\node[paso, below=of d] (e) {\texttt{extraer\_muestras()} y CSV (03a)};
\node[paso, left=of e] (f) {Train / test 75/25 estratificado};
\node[paso, left=of f] (g) {Random Forest y KNN, búsqueda de \texttt{max\_depth} (03b)};
\node[paso, left=of g] (h) {Métricas y mapa del recorte 1};
\node[paso, below=of h] (i) {Validación con recorte 2 y sus muestras (``Yapa'')};
\draw[flecha] (a) -- (b);
\draw[flecha] (b) -- (c);
\draw[flecha] (c) -- (d);
\draw[flecha] (d) -- (e);
\draw[flecha] (e) -- (f);
\draw[flecha] (f) -- (g);
\draw[flecha] (g) -- (h);
\draw[flecha] (h) -- (i);
\end{tikzpicture}
\end{adjustbox}
\caption{Flujo de trabajo seguido.}
\label{fig:flujo}
\end{figure}

\subsection{Lectura de bandas y recortes}
Las bandas se leen directamente del \texttt{.SAFE.zip}, sin descomprimirlo, con la ruta
\texttt{/vsizip/} de GDAL. Se usa una versión de la función \texttt{leer\_bandas\_l2a()} de la Notebook 02c.
Con respecto a la notebook se hicieron tres cambios:
\begin{enumerate}
    \item Las bandas se devuelven siempre en el orden pedido (B02, B03, B04, B08, B11, B12) y no en el orden
    en que aparecen dentro del zip.
    \item El polígono de recorte se reproyecta al CRS de la imagen. En la notebook el CRS estaba escrito a
    mano para Córdoba (UTM 20S), y aquí la escena está en UTM 19N.
    \item B11 y B12 se leen a 20~m y se remuestrean a 10~m de forma bilineal sobre la misma ventana del
    recorte. Así los dos bloques quedan con exactamente el mismo tamaño y se pueden concatenar.
\end{enumerate}
Cada recorte se guardó como un GeoTIFF de seis bandas (\texttt{s2\_recorte1.tif} y
\texttt{s2\_recorte2.tif}), con valor \emph{nodata} igual a 0.

\subsection{Definición de clases y toma de muestras}
Se usaron las seis categorías del primer nivel CIAF. La Tabla~\ref{tab:clases} resume cómo se reconoció
cada una al digitalizar. Se trabajó sobre todo en falso color de vegetación (B08-B04-B03) y se comprobó en
color natural.

\begin{table}[H]
\centering
\caption{Clases utilizadas (nivel 1 CIAF) y criterio de interpretación visual.}
\label{tab:clases}
\small
\renewcommand{\arraystretch}{1.25}
\begin{tabular}{@{}cp{4.2cm}p{9.6cm}@{}}
\toprule
\texttt{Clase} & \texttt{DESC} & Criterio en la imagen \\
\midrule
1 & Construcciones & Casco urbano, vías, aeropuerto. Tonos cian y grises en falso color, textura heterogénea. \\
2 & Cultivos y parcelas & Sobre todo por la forma: lotes geométricos con bordes rectos. En falso color se ven azulados (suelo preparado o con poca cobertura). \\
3 & Vegetación herbácea descubierta & Pastos y sabana. Tonos naranja o pardo rojizo apagados en falso color, textura lisa. \\
4 & Bosques y montes & Rojo intenso en falso color y textura más tupida o rugosa (bosques de galería y relictos). \\
5 & Cuerpos de agua & Río, lagunas, pozos y meandros. Tonos oscuros o azules en falso color, cafés en color natural por los sedimentos. \\
6 & Tierras eriales & Suelo desnudo: playones y bancos de arena del río y parches sin vegetación. \\
\bottomrule
\end{tabular}
\end{table}

En cada recorte se dibujó una capa de polígonos rectangulares (\texttt{muestras\_recorte1.gpkg} y
\texttt{muestras\_recorte2.gpkg}) con los atributos \texttt{Clase} (entero) y \texttt{DESC} (texto). Se
buscó que cada rectángulo quedara dentro de una sola cobertura y lejos de los bordes, para no meter píxeles
mezclados. En el agua esto fue más difícil en los meandros angostos, donde los rectángulos tuvieron que ser
pequeños. Las muestras del recorte 2 se tomaron de forma independiente, sin reutilizar ningún polígono del
recorte 1.

En el recorte 2 había nubes y sombras, pero en el recorte 1 no. Como el modelo solo puede reconocer las
clases con las que fue entrenado, se decidió no muestrearlas.

\begin{figure}[H]
\centering
\includegraphics[width=0.9\linewidth]{muestras_qgis_recorte1.png}
% COLOCAR FIGURA: pantallazo de QGIS con las muestras del recorte 1 sobre la imagen en color natural y la
% tabla de atributos abierta (Clase / DESC), como el que se usó al revisar las 27 primeras muestras.
\caption{Digitalización de las muestras del recorte 1 en QGIS. Los colores siguen la simbología de la Tabla~\ref{tab:clases}.}
\label{fig:qgis}
\end{figure}

% COLOCAR FIGURA (opcional): pantallazo de QGIS con las muestras del recorte 2.

Los polígonos se convirtieron en píxeles con la función \texttt{extraer\_muestras()} de la Notebook 03a.
La función aplica \texttt{mask} con cada polígono y descarta los píxeles de relleno (65535) y los de
\emph{nodata} (0). La Tabla~\ref{tab:muestras} muestra cuántos píxeles quedaron por clase.

\begin{table}[H]
\centering
\caption{Píxeles de muestra por clase en cada recorte.}
\label{tab:muestras}
\small
\begin{tabular}{@{}lrr@{}}
\toprule
Clase & Recorte 1 (entrenamiento) & Recorte 2 (validación) \\
\midrule
Construcciones                  & \num{1814} & \num{92}   \\
Cultivos y parcelas             & \num{708}  & \num{3572} \\
Vegetación herbácea descubierta & \num{1015} & \num{1744} \\
Bosques y montes                & \num{1172} & \num{1527} \\
Cuerpos de agua                 & \num{495}  & \num{313}  \\
Tierras eriales                 & \num{84}   & \num{90}   \\
\midrule
Total                           & \num{5288} & \num{7338} \\
\bottomrule
\end{tabular}
\end{table}

Las proporciones de los dos recortes reflejan lo que hay en cada paisaje. En el recorte 1 domina lo urbano y
en el recorte 2 las parcelas. La clase más débil es tierras eriales, con apenas 84 píxeles en el recorte de
entrenamiento (1,6~\% del total).

\subsection{Análisis de separabilidad}
Antes de entrenar se graficaron las firmas espectrales medias de cada clase (con $\pm$ una desviación
estándar) y un \emph{pairplot} con la dispersión entre pares de bandas.

\begin{figure}[H]
\centering
\includegraphics[width=0.7\linewidth]{firmas_espectrales_recorte1.png}
% COLOCAR FIGURA: firmas espectrales del recorte 1, página 4 del PDF "entregaclasificacion - Colab"
% (salidas/firmas_espectrales_recorte1.png).
\caption{Firmas espectrales medias por clase en el recorte 1.}
\label{fig:firmas}
\end{figure}

\begin{figure}[H]
\centering
\includegraphics[width=0.85\linewidth]{pairplot_recorte1.png}
% COLOCAR FIGURA: pairplot de las seis bandas, páginas 4 y 5 del PDF "entregaclasificacion - Colab"
% (salidas/pairplot_recorte1.png).
\caption{Histogramas y dispersión entre pares de bandas para las muestras del recorte 1.}
\label{fig:pairplot}
\end{figure}

Las firmas se comportan como se esperaba. El agua es la más baja en el NIR y el SWIR. Los bosques tienen el
salto típico entre el rojo y el NIR, y valores bajos en el SWIR. La vegetación herbácea y los cultivos se
parecen bastante entre sí, y por eso se solapan en el espacio de atributos. Construcciones y tierras
eriales tienen reflectancias altas en el visible y el SWIR, con una forma de curva parecida.
% REVISAR contra la Figura de firmas: ajustar esta descripción si alguna clase se ve distinta.

\begin{figure}[H]
\centering
\begin{subfigure}{0.48\linewidth}
    \centering
    \includegraphics[width=\linewidth]{poligonos_recorte1.png}
    \caption{Recorte 1.}
\end{subfigure}\hfill
\begin{subfigure}{0.48\linewidth}
    \centering
    \includegraphics[width=\linewidth]{poligonos_recorte2.png}
    \caption{Recorte 2.}
\end{subfigure}
% COLOCAR FIGURA: polígonos de muestras sobre el falso color, páginas 5 y 6 del PDF
% "entregaclasificacion - Colab" (salidas/poligonos_recorte1.png y poligonos_recorte2.png).
\caption{Polígonos de muestras sobre cada recorte en falso color (B08-B04-B03).}
\label{fig:poligonos}
\end{figure}

\subsection{Entrenamiento y evaluación}
Las muestras del recorte 1 se dividieron en 75~\% para entrenamiento (\num{3966} píxeles) y 25~\% para
prueba (\num{1322}). La división fue estratificada por clase, con mezcla aleatoria y semilla 42, como en la
Notebook 03b. Se entrenaron:
\begin{itemize}
    \item Un \textbf{Random Forest base} con 100 árboles y \texttt{max\_depth=10}.
    \item Una búsqueda de \texttt{max\_depth} entre 2 y 19, con validación cruzada de 5 particiones sobre
    los conjuntos de entrenamiento y de prueba. Con la mejor profundidad se entrenó el \textbf{Random Forest
    final}.
    \item Dos \textbf{KNN} con $k=3$, uno con los datos sin escalar y otro con \texttt{RobustScaler}, para
    comparar.
\end{itemize}
Todos los modelos se evaluaron primero con el 25~\% de prueba del recorte 1. Después se evaluaron con
\emph{todas} las muestras del recorte 2, que el modelo nunca vio. Por último, el Random Forest final se
aplicó a cada recorte completo. Para eso la imagen se reorganiza como una tabla de píxeles por bandas, se
predice y se vuelve a armar la grilla. Los mapas se guardaron como GeoTIFF y como GeoPackage para revisarlos
en QGIS.

% =====================================================================
\section{Resultados}
\label{sec:resultados}
% =====================================================================

\subsection{Profundidad del Random Forest}
La Figura~\ref{fig:profundidad} muestra la exactitud de la validación cruzada para cada profundidad. Hasta
\texttt{max\_depth} $\approx 4$ la mejora es rápida. Después la curva se aplana, y el máximo en el
conjunto de prueba se alcanza en 10 (\num{0.974}), con el mismo valor en 11 y 12. Las curvas de
entrenamiento y de prueba van casi juntas, con una diferencia de menos de un punto, así que dentro del
recorte 1 no se ve sobreajuste. Por eso el Random Forest final quedó con \texttt{max\_depth=10}, igual que
el modelo base, y los dos dan exactamente las mismas métricas.

\begin{figure}[H]
\centering
\begin{tikzpicture}
\begin{axis}[
    width=0.72\linewidth, height=6cm,
    xlabel={Profundidad máxima (\texttt{max\_depth})}, ylabel={Exactitud (CV, 5 particiones)},
    xmin=1.5, xmax=19.5, ymin=0.90, ymax=0.985,
    xtick={2,4,...,18},
    yticklabel style={/pgf/number format/fixed, /pgf/number format/precision=2,
                      /pgf/number format/use comma},
    grid=major, grid style={lightGray},
    legend pos=south east, legend style={font=\footnotesize},
]
\addplot[mainBlue, thick, mark=*, mark size=1.4pt] coordinates {
(2,0.914)(3,0.943)(4,0.964)(5,0.964)(6,0.966)(7,0.968)(8,0.970)(9,0.973)(10,0.974)
(11,0.974)(12,0.974)(13,0.972)(14,0.973)(15,0.973)(16,0.973)(17,0.973)(18,0.973)(19,0.973)};
\addlegendentry{Prueba}
\addplot[orange!80!black, thick, dashed, mark=square*, mark size=1.2pt] coordinates {
(2,0.909)(3,0.940)(4,0.961)(5,0.969)(6,0.972)(7,0.974)(8,0.975)(9,0.976)(10,0.976)
(11,0.977)(12,0.977)(13,0.977)(14,0.978)(15,0.977)(16,0.977)(17,0.977)(18,0.978)(19,0.978)};
\addlegendentry{Entrenamiento}
\draw[red, dashed] (axis cs:10,0.90) -- (axis cs:10,0.985);
\end{axis}
\end{tikzpicture}
% Figura generada con los valores impresos por el script (página 7 del PDF "entregaclasificacion - Colab").
% Si se prefiere la del script, usar salidas/rf_exactitud_vs_profundidad.png (páginas 7-8 del PDF).
\caption{Exactitud de la validación cruzada del Random Forest en función de la profundidad máxima. La línea roja marca la profundidad elegida.}
\label{fig:profundidad}
\end{figure}

\subsection{Evaluación sobre el recorte 1}
En el conjunto de prueba del recorte 1, los tres modelos dan resultados muy altos y casi iguales
(Tabla~\ref{tab:global}). En el Random Forest, construcciones y cuerpos de agua quedan perfectos
(F1 = 1,00) y cultivos y parcelas llega a 0,98. Las confusiones que quedan están entre vegetación herbácea
y bosques (F1 de 0,95 y 0,96) y en tierras eriales, que tiene un \emph{recall} de 0,86: de 21 píxeles de
prueba, tres se fueron a otras clases. KNN acierta lo mismo en la exactitud global, pero con tierras
eriales le va peor (\emph{recall} 0,71).

\begin{table}[H]
\centering
\caption{Métricas por clase del Random Forest (\texttt{max\_depth=10}) en el conjunto de prueba del recorte 1.}
\label{tab:rf_test}
\small
\begin{tabular}{@{}lcccr@{}}
\toprule
Clase & Precisión & Recall & F1 & Soporte \\
\midrule
Construcciones                  & 1,00 & 1,00 & 1,00 & 453 \\
Cultivos y parcelas             & 0,99 & 0,98 & 0,98 & 177 \\
Vegetación herbácea descubierta & 0,95 & 0,95 & 0,95 & 254 \\
Bosques y montes                & 0,96 & 0,97 & 0,96 & 293 \\
Cuerpos de agua                 & 1,00 & 1,00 & 1,00 & 124 \\
Tierras eriales                 & 1,00 & 0,86 & 0,92 & 21  \\
\midrule
Promedio macro                  & 0,98 & 0,96 & 0,97 & \num{1322} \\
\bottomrule
\end{tabular}
\end{table}

\begin{figure}[H]
\centering
\includegraphics[width=0.62\linewidth]{confusion_rf_mejor_test.png}
% COLOCAR FIGURA: matriz de confusión del RF max_depth=10 en el test del recorte 1, página 8 del PDF
% "entregaclasificacion - Colab" (salidas/confusion_rf_mejor_test.png).
% Las matrices de KNN (páginas 8-10) pueden ir en un anexo si se quiere.
\caption{Matriz de confusión del Random Forest en el conjunto de prueba del recorte 1.}
\label{fig:conf_test}
\end{figure}

\subsection{Validación con el recorte 2}
Con las muestras independientes del recorte 2, la exactitud del Random Forest baja de \num{0.979} a
\num{0.668}, y el kappa de \num{0.972} a \num{0.531} (Tabla~\ref{tab:global}). KNN sin escalar queda un
poco por debajo, con \num{0.642} y $\kappa=\num{0.503}$.

\begin{table}[H]
\centering
\caption{Exactitud global e índice kappa de cada modelo.}
\label{tab:global}
\small
\begin{tabular}{@{}lcccc@{}}
\toprule
 & \multicolumn{2}{c}{Prueba recorte 1} & \multicolumn{2}{c}{Validación recorte 2} \\
\cmidrule(lr){2-3}\cmidrule(lr){4-5}
Modelo & Exactitud & Kappa & Exactitud & Kappa \\
\midrule
Random Forest (\texttt{max\_depth=10}) & \num{0.979} & \num{0.972} & \num{0.668} & \num{0.531} \\
KNN ($k=3$) sin escalar                & \num{0.979} & \num{0.972} & \num{0.642} & \num{0.503} \\
KNN ($k=3$) con \texttt{RobustScaler}  & \num{0.977} & \num{0.969} & --- & --- \\
\bottomrule
\end{tabular}
% COMPLETAR: los valores de KNN escalado en la validación no salieron en el PDF exportado (la página 14
% solo tiene la matriz). Copiarlos de la salida del script en Colab.
\end{table}

\begin{table}[H]
\centering
\caption{Métricas por clase del Random Forest con las muestras de validación del recorte 2.}
\label{tab:rf_val}
\small
\begin{tabular}{@{}lcccr@{}}
\toprule
Clase & Precisión & Recall & F1 & Soporte \\
\midrule
Construcciones                  & 0,21 & 1,00 & 0,35 & 92   \\
Cultivos y parcelas             & 0,81 & 0,55 & 0,66 & \num{3572} \\
Vegetación herbácea descubierta & 0,48 & 0,72 & 0,58 & \num{1744} \\
Bosques y montes                & 0,93 & 0,90 & 0,92 & \num{1527} \\
Cuerpos de agua                 & 0,75 & 0,63 & 0,68 & 313  \\
Tierras eriales                 & 0,01 & 0,01 & 0,01 & 90   \\
\midrule
Promedio macro                  & 0,53 & 0,64 & 0,53 & \num{7338} \\
\bottomrule
\end{tabular}
\end{table}

\begin{figure}[H]
\centering
\includegraphics[width=0.62\linewidth]{confusion_rf_mejor_validacion.png}
% COLOCAR FIGURA: matriz de confusión del RF en la validación con el recorte 2, página 12 del PDF
% "entregaclasificacion - Colab" (salidas/confusion_rf_mejor_validacion.png).
\caption{Matriz de confusión del Random Forest con las muestras del recorte 2.}
\label{fig:conf_val}
\end{figure}

Bosques y montes es la única clase que se mantiene firme, con F1 de 0,92. Hay tres clases con problemas
claros:
\begin{itemize}
    \item \textbf{Construcciones:} todas sus muestras se reconocen (\emph{recall} 1,00), pero la precisión
    cae a 0,21. Esto quiere decir que, de cada cinco píxeles que el modelo llamó construcción, solo uno lo
    era. Las métricas indican que unos 340 píxeles de otras clases se fueron a construcción.
    \item \textbf{Tierras eriales:} prácticamente no se reconoce (F1 = 0,01).
    \item \textbf{Cultivos y vegetación herbácea:} cultivos pierde el 45~\% de sus píxeles
    (\emph{recall} 0,55), y la herbácea tiene precisión de 0,48. Las cuentas cuadran con que buena parte de
    los cultivos terminó clasificada como herbácea.
\end{itemize}
% VERIFICAR en la matriz (pág. 12) de qué clases vienen los falsos positivos de construcción
% (se espera que sean sobre todo tierras eriales y cultivos/parcelas).

\subsection{Mapas clasificados}

\begin{figure}[H]
\centering
\includegraphics[width=0.75\linewidth]{mapa_recorte1_rf.png}
% COLOCAR FIGURA: mapa clasificado del recorte 1 con RF, páginas 10 y 11 del PDF
% "entregaclasificacion - Colab" (salidas/mapa_recorte1_rf.png).
\caption{Mapa clasificado del recorte 1 (Random Forest, \texttt{max\_depth=10}).}
\label{fig:mapa1}
\end{figure}

\begin{figure}[H]
\centering
\includegraphics[width=0.75\linewidth]{mapa_recorte2_rf.png}
% COLOCAR FIGURA: mapa clasificado del recorte 2 con RF. No quedó en el PDF exportado (se corta en la
% página 14); tomarlo de salidas/mapa_recorte2_rf.png o hacer el pantallazo en QGIS con
% clasificacion_recorte2_rf.gpkg.
\caption{Mapa clasificado del recorte 2 (Random Forest, \texttt{max\_depth=10}).}
\label{fig:mapa2}
\end{figure}

% =====================================================================
\section{Análisis de resultados}
\label{sec:analisis}
% =====================================================================

\subsection{Lo que el modelo clasificó bien}
En los dos recortes, \textbf{los bosques y la vegetación herbácea descubierta} se ven bien delimitados en
los mapas. En el caso de los bosques coincide con las métricas: es la clase más estable al pasar de un
recorte al otro. Tiene sentido, porque su combinación de reflectancia alta en el NIR y baja en el SWIR es
muy distinta de la de las demás clases.

En el recorte 1, el modelo reconoció como agua \textbf{las lagunas alargadas que acompañan al río}
(Figura~\ref{fig:lagunas}). No se daba por hecho, porque no todas tenían muestras y algunas tienen
vegetación en los bordes. También quedaron bien \textbf{el casco urbano y el aeropuerto}
(Figura~\ref{fig:urbano}): la pista y las construcciones salen como una sola mancha de construcciones,
cortada por los corredores de vegetación.

\begin{figure}[H]
\centering
\includegraphics[width=0.45\linewidth]{lagunas_recorte1.png}
% COLOCAR FIGURA: pantallazo de QGIS del mapa del recorte 1 en la zona de las lagunas alargadas paralelas
% al río (clasificadas en azul como cuerpos de agua).
\caption{Lagunas del recorte 1 clasificadas correctamente como cuerpos de agua.}
\label{fig:lagunas}
\end{figure}

\begin{figure}[H]
\centering
\includegraphics[width=0.9\linewidth]{urbano_aeropuerto_recorte1.png}
% COLOCAR FIGURA: pantallazo doble del aeropuerto y la zona urbana (mapa clasificado a la izquierda,
% color natural a la derecha).
\caption{Zona urbana y aeropuerto en el recorte 1: clasificación (izquierda) e imagen en color natural (derecha).}
\label{fig:urbano}
\end{figure}

\subsection{Suelos descubiertos clasificados como construcción}
El error más visible está en las superficies con el suelo descubierto: los playones y bancos de arena del
río y algunas parcelas sin cobertura. El modelo las clasifica como \textbf{construcciones}, a veces con un
borde de tierras eriales alrededor (Figuras~\ref{fig:playones} y~\ref{fig:parche}).

\begin{figure}[H]
\centering
\begin{subfigure}{0.48\linewidth}
    \centering
    \includegraphics[width=\linewidth]{playones_clasificacion.png}
    \caption{Clasificación.}
\end{subfigure}\hfill
\begin{subfigure}{0.48\linewidth}
    \centering
    \includegraphics[width=\linewidth]{playones_color_natural.png}
    \caption{Color natural.}
\end{subfigure}
% COLOCAR FIGURA: pantallazos del meandro con los playones de arena (clasificados en rojo) y la misma zona
% en color natural, con el rectángulo de la muestra de tierras eriales.
\caption{Playones del río clasificados como construcciones. El rectángulo es una muestra de tierras eriales.}
\label{fig:playones}
\end{figure}

\begin{figure}[H]
\centering
\begin{subfigure}{0.48\linewidth}
    \centering
    \includegraphics[width=\linewidth]{parche_suelo_clasificacion.png}
    \caption{Clasificación.}
\end{subfigure}\hfill
\begin{subfigure}{0.48\linewidth}
    \centering
    \includegraphics[width=\linewidth]{parche_suelo_color_natural.png}
    \caption{Color natural.}
\end{subfigure}
% COLOCAR FIGURA: pantallazos del parche de suelo descubierto junto a la vía (rojo con borde café en la
% clasificación) y la misma zona en color natural, con la muestra de tierras eriales.
\caption{Parche de suelo descubierto clasificado como construcciones con borde de tierras eriales.}
\label{fig:parche}
\end{figure}

Hay varias razones que se suman:
\begin{enumerate}
    \item \textbf{Firmas parecidas.} La arena seca y el suelo arado reflejan mucho en el visible y en el
    SWIR, igual que los techos, el pavimento y la pista del aeropuerto. En el espacio de seis bandas,
    esas superficies quedan muy cerca de las construcciones.
    \item \textbf{Muy pocas muestras de tierras eriales.} En el recorte 1 solo hubo un polígono de esta clase
    (84 píxeles, contra \num{1814} de construcciones). Los árboles del Random Forest casi no ``vieron''
    suelo desnudo, así que, ante un píxel brillante, la clase que más votos recibe es construcción. El borde
    café que aparece alrededor de los playones corresponde a los píxeles de transición, más oscuros o más
    húmedos, que se parecen a la única muestra de eriales.
    \item \textbf{La muestra de eriales no representa toda la variabilidad.} Un parche de suelo en medio de
    un potrero no responde igual que un banco de arena seca del río. Por eso, en la validación, tierras
    eriales queda en F1 = 0,01.
\end{enumerate}

En el recorte 1 se ve el mismo problema en una zona de parcelas (Figura~\ref{fig:parcela1}). Un lote con el
suelo descubierto quedó como construcción y tierras eriales. El rectángulo amarillo, una muestra de cultivos
y parcelas que está al lado, sí quedó bien clasificado.

\begin{figure}[H]
\centering
\begin{subfigure}{0.48\linewidth}
    \centering
    \includegraphics[width=\linewidth]{parcela_recorte1_color_natural.png}
    \caption{Color natural.}
\end{subfigure}\hfill
\begin{subfigure}{0.48\linewidth}
    \centering
    \includegraphics[width=\linewidth]{parcela_recorte1_clasificacion.png}
    \caption{Clasificación.}
\end{subfigure}
% COLOCAR FIGURA: pantallazos de la parcela del recorte 1 con la muestra amarilla de cultivos (color natural
% y clasificación, donde aparece la mancha roja con borde café).
\caption{Parcela con suelo descubierto en el recorte 1 clasificada como construcciones.}
\label{fig:parcela1}
\end{figure}

\subsection{El ``invernadero'' del recorte 2}
En el recorte 2 se tomó como construcción una zona que en falso color se veía cian brillante, con una forma
regular, y que parecía un invernadero (Figura~\ref{fig:invernadero}). Revisándola en color natural, se ve
más bien como un lote de color café claro, con textura de suelo y sin techos. Lo más probable es que sea
\textbf{una parcela con el suelo descubierto}, arada o preparada para siembra, o recién cosechada. En CIAF
correspondería a cultivos y parcelas (temporales, en barbecho) o, si está abandonada, a tierras eriales. Sin
verificación en campo no se puede confirmar.

\begin{figure}[H]
\centering
\begin{subfigure}{0.40\linewidth}
    \centering
    \includegraphics[width=\linewidth]{invernadero_color_natural.png}
    \caption{Color natural.}
\end{subfigure}\hspace{0.05\linewidth}
\begin{subfigure}{0.40\linewidth}
    \centering
    \includegraphics[width=\linewidth]{invernadero_falso_color.png}
    \caption{Falso color (B08-B04-B03).}
\end{subfigure}
% COLOCAR FIGURA: pantallazos de la muestra roja del supuesto invernadero en color natural y en falso color.
\caption{Zona del recorte 2 que se muestreó como construcción (invernadero) y que probablemente es una parcela con suelo descubierto.}
\label{fig:invernadero}
\end{figure}

Este error de etiquetado ayuda a entender las métricas de construcciones en la validación. De las 92
muestras de esa clase en el recorte 2, una parte corresponde a esa parcela. El modelo la clasificó como
construcción, que era la etiqueta que se le había puesto. Así, el \emph{recall} de 1,00 está inflado: el
modelo y la etiqueta se equivocan en el mismo sentido. Al mismo tiempo, la baja precisión (0,21) muestra que
el modelo manda a construcción muchas otras superficies brillantes.

\subsection{Nubes y sombras}
Las nubes del recorte 2 salieron clasificadas como construcciones (Figura~\ref{fig:nubes}). Era de esperar:
una nube refleja mucho en todas las bandas, y entre las seis clases entrenadas la más brillante es
construcción. Como no había nubes ni sombras en el recorte 1, ninguna clase las representa, y el modelo las
asigna a la clase más parecida que conoce. Este error no aparece en las métricas, porque no se tomaron
muestras sobre nubes, pero sí ensucia el mapa. Las sombras, por ser oscuras, pueden terminar como agua o
como bosque.
% REVISAR en el mapa del recorte 2 cómo quedaron las sombras y ajustar la última frase.

\begin{figure}[H]
\centering
\includegraphics[width=0.38\linewidth]{nubes_recorte2.png}
% COLOCAR FIGURA: pantallazo del mapa del recorte 2 donde las nubes aparecen en rojo (construcciones).
\caption{Nubes del recorte 2 clasificadas como construcciones.}
\label{fig:nubes}
\end{figure}

\subsection{Cultivos y vegetación herbácea}
La otra gran fuente de error es la confusión entre \textbf{cultivos y parcelas} y \textbf{vegetación
herbácea descubierta}. Se veía venir desde la toma de muestras. Los cultivos se identificaron más que todo
por la forma de las parcelas y por los tonos azulados del suelo preparado, y la herbácea por los tonos
naranja. Pero un cultivo ya crecido tiene una firma muy parecida a la de un pasto. El clasificador trabaja
píxel a píxel y no ``ve'' la forma del lote, así que pierde justamente el criterio que se usó para
diferenciarlas.

A esto se suma que el recorte 1 tenía pocas parcelas (708 píxeles) y el recorte 2 está dominado por ellas
(\num{3572}). El modelo aprendió los cultivos con pocos ejemplos y casi todos del mismo tipo. Cuando
encontró parcelas en otros estados en el recorte 2, muchas las clasificó como herbácea.

\subsection{Cuerpos de agua}
El río se clasifica sin problema, pero en el recorte 2 el \emph{recall} del agua baja a 0,63. Coincide con
lo que se notó al muestrear: el agua es fácil de ver en el cauce principal, pero no tanto en algunas
lagunas, pozos y meandros. Ahí el agua es poco profunda, tiene sedimentos o vegetación, y los píxeles de
borde se mezclan con la orilla. Como las muestras de los meandros fueron pequeñas, una parte de sus píxeles
probablemente ya incluía orilla.

\subsection{Diferencia entre las dos validaciones}
La caída de 0,98 a 0,67 en la exactitud no significa que la primera evaluación esté mal hecha. Mide cosas
distintas. En el recorte 1, los píxeles de prueba salen de los mismos polígonos que los de entrenamiento y
son casi idénticos a ellos, por lo que el resultado es optimista. El recorte 2 tiene otro paisaje, otras
proporciones de clases, nubes y un error de etiquetado, y muestra mejor cómo se comportaría el modelo en una
zona nueva. Por eso tenía sentido hacer la segunda validación que pide la consigna.

\begin{analisisBox}{Segunda clasificación (en preparación)}
A partir de lo anterior, se va a realizar una segunda clasificación con estos cambios:
\begin{itemize}
    \item Corregir la muestra del supuesto invernadero del recorte 2, reclasificándola como cultivos y
    parcelas (o eliminándola si no se puede confirmar).
    \item Aumentar las muestras de \textbf{tierras eriales} en el recorte 1, incluyendo playones y bancos de
    arena del río y parcelas con el suelo descubierto, para que el modelo aprenda a separarlas de las
    construcciones.
    \item Tomar más muestras de \textbf{cultivos y parcelas} en el recorte 1, en distintos estados (con
    cultivo, arado, en barbecho). Si hace falta, se moverá o ampliará el recorte para incluir más lotes.
    \item Tratar las nubes y sombras, ya sea con una clase propia muestreada en ambos recortes o
    enmascarándolas con la capa de clasificación de escena (SCL) que trae el producto L2A.
    \item Evaluar si agregar índices como el NDVI y el NDWI como atributos mejora la separación entre
    agua, vegetación y suelo.
\end{itemize}
\end{analisisBox}

% =====================================================================
\section{Conclusiones}
% =====================================================================

\begin{itemize}
    \item Fue posible reproducir con una escena propia, sobre Arauca, el flujo de trabajo del curso:
    lectura de Sentinel-2 L2A desde el zip, recortes, muestras, entrenamiento, evaluación y mapas. Solo hubo
    que adaptar el sistema de coordenadas y el orden de las bandas.
    \item Random Forest y KNN dieron prácticamente lo mismo en el recorte de entrenamiento (exactitud cercana
    a 0,98). En la validación externa, el Random Forest fue un poco mejor (0,668 contra 0,642). La
    profundidad óptima quedó en 10, sin señales de sobreajuste.
    \item La validación con un recorte distinto fue clave. Las métricas internas daban casi un modelo
    perfecto, mientras que con muestras independientes la exactitud bajó a 0,67 y el kappa a 0,53. Ahí se
    vieron los problemas reales del modelo.
    \item Las clases más confiables fueron bosques y montes y, en el recorte 1, construcciones y cuerpos de
    agua. Los errores principales vienen de los suelos descubiertos y brillantes clasificados como
    construcción, de la confusión entre cultivos y vegetación herbácea, y de las nubes, que no estaban en
    el entrenamiento.
    \item Buena parte de esos errores se origina en las muestras más que en el algoritmo: pocos píxeles de
    tierras eriales, pocas parcelas en el recorte de entrenamiento y una muestra mal interpretada. Por eso
    la segunda clasificación se enfoca en mejorar las muestras.
    \item Al trabajar píxel a píxel se pierde información que el intérprete sí usa, como la forma de las
    parcelas o la textura del bosque. Ese es un límite de este enfoque para separar cultivos de pastos.
\end{itemize}

% =====================================================================
\begin{thebibliography}{9}
% =====================================================================
\bibitem[Breiman(2001)]{breiman2001}
Breiman, L. (2001). Random forests. \textit{Machine Learning}, 45(1), 5--32.

\bibitem[{CIAF}(s.f.)]{ciaf}
Centro de Investigación en Percepción Remota (CIAF) (s.f.). \textit{Categorías de cobertura y uso de la
tierra por nivel de detalle}. Instituto Geográfico Agustín Codazzi. Material de clase, Universidad
Distrital Francisco José de Caldas.
% COMPLETAR: referencia exacta del documento CIAF que les dieron en clase.

\bibitem[Chuvieco(2010)]{chuvieco2010}
Chuvieco, E. (2010). \textit{Teledetección ambiental: la observación de la Tierra desde el espacio}
(3.ª ed.). Ariel.

\bibitem[Cohen(1960)]{cohen1960}
Cohen, J. (1960). A coefficient of agreement for nominal scales. \textit{Educational and Psychological
Measurement}, 20(1), 37--46.

\bibitem[Congalton(1991)]{congalton1991}
Congalton, R. G. (1991). A review of assessing the accuracy of classifications of remotely sensed data.
\textit{Remote Sensing of Environment}, 37(1), 35--46.

\bibitem[{Copernicus}(2025)]{copernicus}
Copernicus Data Space Ecosystem (2025). \textit{Copernicus Browser}.
\url{https://browser.dataspace.copernicus.eu/}

\bibitem[Drusch et~al.(2012)]{drusch2012}
Drusch, M., Del Bello, U., Carlier, S., Colin, O., Fernandez, V., Gascon, F., et~al. (2012). Sentinel-2:
ESA's optical high-resolution mission for GMES operational services. \textit{Remote Sensing of
Environment}, 120, 25--36.

\bibitem[{ESA}(2015)]{esa2015}
European Space Agency (2015). \textit{Sentinel-2 User Handbook}. ESA Standard Document.

\bibitem[{Instituto Gulich}(s.f.)]{gulich}
Instituto de Altos Estudios Espaciales Mario Gulich (s.f.). \textit{Procesamiento digital de imágenes
satelitales y datos geoespaciales con Python}, Bloque 03: presentaciones de clase y Notebooks 02c, 03a y
03b. Campus Virtual, Universidad Nacional de Córdoba -- CONAE.

\bibitem[Pedregosa et~al.(2011)]{pedregosa2011}
Pedregosa, F., Varoquaux, G., Gramfort, A., Michel, V., Thirion, B., Grisel, O., et~al. (2011).
Scikit-learn: Machine learning in Python. \textit{Journal of Machine Learning Research}, 12, 2825--2830.
\end{thebibliography}

\appendix

% =====================================================================
\section{Organización del código fuente}
\label{anx:codigo}
% =====================================================================
El código completo está en \texttt{entregaclasificacion.py}, en el repositorio
\url{https://github.com/AlexandraUD/ClasificacionSupervisada} (Anexo~\ref{anx:archivos}), y no se transcribe
en este documento. Usa \texttt{rasterio}, \texttt{geopandas}, \texttt{scikit-learn} \citep{pedregosa2011},
\texttt{matplotlib} y \texttt{seaborn}. Sus funciones se agrupan así:
\begin{itemize}
    \item \textbf{Lectura de Sentinel-2 (02c):} \texttt{leer\_bandas\_l2a}, \texttt{leer\_bandas\_remuestreadas},
    \texttt{generar\_recorte}.
    \item \textbf{Muestras (03a):} \texttt{extraer\_muestras}, \texttt{leer\_vector\_muestras},
    \texttt{graficar\_firmas\_espectrales}, \texttt{graficar\_recorte\_con\_poligonos}.
    \item \textbf{Clasificación y evaluación (03b):} \texttt{evaluar} (matriz de confusión, reporte, exactitud
    y kappa) y \texttt{clasificar\_imagen} (mapa en GeoTIFF y GeoPackage).
    \item \textbf{Visualización y utilidades:} \texttt{escalar\_imagen}, \texttt{escalar\_imagen\_multibanda},
    \texttt{asignar\_colores} (colores de la simbología de QGIS), \texttt{preparar\_zip}, \texttt{buscar\_vector}.
\end{itemize}
En Colab se ejecuta con:
\begin{lstlisting}[language=Python]
!pip install -q rasterio geopandas scikit-learn scipy matplotlib seaborn gdown
!wget -q -O entregaclasificacion.py https://raw.githubusercontent.com/AlexandraUD/ClasificacionSupervisada/main/entregaclasificacion.py
%run entregaclasificacion.py
\end{lstlisting}

% =====================================================================
\section{Archivos de la entrega}
\label{anx:archivos}
% =====================================================================
\begin{table}[H]
\centering
\caption{Archivos que acompañan al informe.}
\small
\renewcommand{\arraystretch}{1.2}
\begin{tabular}{@{}p{5.6cm}p{9.8cm}@{}}
\toprule
Archivo & Contenido \\
\midrule
\texttt{entregaclasificacion.py} & Script completo: lectura, recortes, muestras, entrenamiento, evaluación y mapas. \\
\texttt{recorte1.gpkg}, \texttt{recorte2.gpkg} & Polígonos de los dos recortes de $5\times5$~km. \\
\texttt{muestras\_recorte1.gpkg} & Muestras de entrenamiento (campos \texttt{Clase} y \texttt{DESC}). \\
\texttt{muestras\_recorte2.gpkg} & Muestras independientes de validación. \\
\texttt{simblogia.qml} & Simbología de QGIS de las clases. \\
\texttt{salidas/s2\_recorte1.tif}, \texttt{s2\_recorte2.tif} & Recortes de seis bandas a 10~m. \\
\texttt{salidas/muestras\_recorte*.csv} & Píxeles de muestra extraídos. \\
\texttt{salidas/clasificacion\_recorte*\_rf.tif / .gpkg} & Mapas clasificados (raster y vector). \\
\texttt{salidas/*.png} & Figuras generadas por el script. \\
\texttt{entregaclasificacion - Colab.pdf} & Salida completa de la ejecución en Colab. \\
\bottomrule
\end{tabular}
\end{table}

\end{document}
